\section{Evaluation}
\label{sec:evaluation}

\subsection{Research Questions}

The evaluation is organized around three research questions. Each RQ is answered
only after a source oracle defines the workload behavior, \namei and FUSE rows
run the same correctness condition, and the boundary account records which
filesystem responsibilities each mechanism owns.
The unit of evidence is a complete source-oracle matrix, not a catalog of
one-off baselines or mechanism smoke tests.

\begin{description}
\item[RQ1] Can a narrow VFS name-resolution extension express real
state-dependent path-view policies without taking over filesystem semantics?
\item[RQ2] What is the cost of putting programmable policy on the VFS
name-resolution path compared with a feature-equivalent FUSE policy
implementation?
\item[RQ3] Does \namei provide a narrower verifier-bounded, fail-closed
ownership boundary than building a custom or stackable filesystem when the
needed policy is only name resolution?
\end{description}

\subsection{Experimental Setup}

All \namei mechanism runs execute in KVM~\cite{linux_kvm} with the modified
kernel and the real \code{cgroup/namei_ext} attach path. Make targets own
workload execution, comparison construction, result collection, and report
generation. Each run records the machine, kernel, image, policy, workload
inputs, stdout/stderr, dmesg, per-operation lookup/readdir events, and
comparison-specific setup observations.
Operation-weighted metrics weight each policy action by its observed frequency
in the workload trace.

Correctness gates every row. A source oracle is the correctness condition
extracted from the original source system. A same-oracle mechanism comparison
requires the executed \namei and feature-equivalent FUSE rows to pass that same
condition, not weaker mechanism-specific tests. RQ3 boundary accounts use the
same source oracle to compare responsibilities and containment rather than to
claim an executed custom/stackable correctness row. Reproduced or replayed
source oracles can become evaluation rows. Source inspection and paper-derived
evidence remain context unless they are backed by a real KVM oracle.

\paragraph{Measurement protocol.}
For each final reported run, the artifact records hardware and kernel
configuration, kernel commit, KVM image digest, filesystem configuration, run
count, warmup rule, confidence-interval method, cache-hot/cache-cold protocol,
feature-equivalent FUSE implementation, and custom or stackable filesystem
boundary account for each reportable workload.

\subsection{RQ1: Expressiveness And Sufficiency}

RQ1 asks whether the VFS boundary can express path-view transitions extracted
from source oracles while preserving lower-filesystem semantics. A reportable
workload must run the complete source lifecycle or fixed replay needed by its
oracle, while \namei evidence is limited to lookup and readdir effects.
Runtime orchestration, storage state, synthetic contents, write handling, and
source-system lifecycle behavior remain delegated to the source system or
lower filesystem.

\begin{table}[t]
\centering
\scriptsize
\begin{tabular}{@{}L{0.18\linewidth}L{0.25\linewidth}L{0.22\linewidth}L{0.25\linewidth}@{}}
\hline
Workload & Source oracle and path-view effect & Lower-filesystem preservation check & RQ1 result slot \\
\hline
Agent workspace lifecycle~\cite{agentfs_repo} & AgentFS-derived branch/checkpoint
state selects or hides branch-visible objects during lookup and readdir. The
oracle checks the final tree, whiteout/hide effects, symlink metadata,
lookup/readdir coherence, and branch-visible state. & Lower filesystem keeps
permissions, writes, symlink metadata, file data, and directory operations
outside the policy. & \textbf{Result slot:} reviewed KVM
pass/fail matrix, operation-weighted lookup/readdir trace, and lower-filesystem
preservation rows. \\
Environment/cache transition~\cite{swe_factory_repo,menvagent_repo,swe_rebench_v2_repo} & Cache epoch selects verified local or canonical fallback objects, or hides stale/corrupt local objects. The oracle checks the source build/test result and stale/corrupt non-use. & Source evaluator and lower filesystem keep cache population, content validation, build/test execution, writes, and data path. & \textbf{Result slot:} KVM cache-state matrix after final-file target support, with hit, miss, stale, corrupt, and epoch-update rows. \\
Service/config transition~\cite{kubernetes_projected_volumes,kubernetes_configmaps,kubernetes_secrets} & Included only if a concrete source oracle depends on lookup-time object selection at a stable service path. & Reload protocol, generated contents, and data path remain outside the hook. & \textbf{Conditional:} no RQ1 row unless a source oracle is admitted. \\
\hline
\end{tabular}
\caption{RQ1 evidence slots. A workload becomes evidence only after the
source-oracle row passes through the real KVM attach path and lower-filesystem
semantics are preserved.}
\label{tab:rq1-evidence}
\end{table}

RQ1 is answered positively when the admitted Agent workspace and
environment/cache rows pass their source oracles through the real KVM attach
path while preserving lookup/readdir coherence, lower-filesystem behavior, and
operation-weighted traces. RQ1 fails for a workload if the oracle requires
synthetic contents, data-path mediation, write-conflict handling, or metadata
ownership that cannot be expressed as bounded name-resolution policy. Result
slot reports the pass/fail outcome, trace evidence, and lower-filesystem
preservation cells in
Table~\ref{tab:rq1-evidence}.
The environment/cache matrix remains a primary target of the hypothesis, but it
is not reported as evidence until final-file target selection is implemented and
reviewed.

\subsection{RQ2: Cost Versus FUSE}

RQ2 compares \namei with a feature-equivalent FUSE policy service under the same
source oracle, isolating the request-path cost of moving the same decision out
of VFS name resolution. For each final reported workload, the FUSE row must
implement the same source policy and pass the same oracle before any overhead
is interpreted. Non-name-resolution lifecycle
behavior is delegated identically in both rows, so the comparison isolates the
placement of the lookup/readdir policy rather than giving either mechanism a
different runtime.

\begin{table}[t]
\centering
\scriptsize
\begin{tabular}{@{}L{0.2\linewidth}L{0.22\linewidth}L{0.22\linewidth}L{0.26\linewidth}@{}}
\hline
Workload & Compared rows & Metrics & RQ2 result slot \\
\hline
Agent workspace lifecycle & no-hook/lower-filesystem control, \namei policy in KVM,
feature-equivalent FUSE policy service & lookup, open, stat, access, exec, and
readdir latency, macro lifecycle runtime, action mix, and FUSE request-path
cost & \textbf{Result slot:} correctness-gated latency and macro table with
confidence intervals and FUSE option record. \\
Environment/cache transition & no-hook/lower-filesystem control, \namei cache policy in
KVM, feature-equivalent FUSE cache-policy service & cache-hot and cache-cold
lookup, open, stat, access, and readdir latency, build/test macro runtime, and
operation-weighted policy invocation rate & \textbf{Result slot:}
correctness-gated cache-state overhead table after final-file target support. \\
\hline
\end{tabular}
\caption{RQ2 cost slots. FUSE numbers are interpreted only after the FUSE row
implements the same source policy and passes the same source oracle.}
\label{tab:rq2-evidence}
\end{table}

The reported metrics are pass-through and action-specific \namei overhead,
lookup/open/stat/access/exec/readdir latency, cache-hot and cache-cold behavior,
macro workload runtime, operation-weighted invocation rate, matched action mix,
and FUSE daemon/request-path cost. No-hook and pass-through controls separate
fixed hook cost from policy cost.

RQ2 supports the paper's boundary claim if the same-oracle \namei
rows have lower lookup/readdir placement cost than feature-equivalent FUSE
without losing correctness or changing the delegated runtime. RQ2 weakens that
claim if FUSE passes the same oracle with comparable or lower cost under fair
caching and passthrough settings. If a feature-equivalent FUSE row cannot be
built, that workload is not interpretable for RQ2 and is reported as a
methodology gap, not as overhead evidence. The corresponding result slot reports
the latency,
macro-runtime, action-mix, and FUSE-option cells in
Table~\ref{tab:rq2-evidence}.

\subsection{RQ3: Safety And Boundary Versus Custom Or Stackable Filesystems}

RQ3 treats custom and stackable filesystems as ownership-boundary comparisons.
For the same source oracle, the question is which filesystem responsibilities
the developer must implement or operate, and whether \namei can keep those
responsibilities in the VFS, lower filesystem, or source runtime. Policy-output
safety and containment mean that the eBPF verifier bounds policy execution and
the kernel validates names, action types, target identifiers, and selection
scope. Malformed decisions fail closed, while lower filesystems still enforce
permissions, data operations, writes, and persistence. Source systems define the
required responsibilities, and the custom/stackable boundary account asks which
of those responsibilities would move into a filesystem boundary.

\begin{table}[t]
\centering
\scriptsize
\begin{tabular}{@{}L{0.18\linewidth}L{0.25\linewidth}L{0.25\linewidth}L{0.22\linewidth}@{}}
\hline
Evidence item & \namei boundary & Custom/stackable boundary & RQ3 result slot \\
\hline
Same-oracle path-view policy & Verified eBPF returns bounded lookup/readdir
actions, and the kernel validates outputs and resumes ordinary VFS execution. &
Filesystem boundary owns or interposes on methods such as lookup, readdir,
create, unlink, rename, open/read/write, or metadata state when the source
oracle requires them. &
\textbf{Result slot:} workload-specific ownership table after RQ1 oracle
passes. \\
Malformed decision containment & Unsupported actions, invalid targets, and
malformed names are rejected on the real \code{cgroup/namei_ext} path. & Custom or
stackable implementations must provide their own validation and failure
behavior across the methods they own. & \textbf{Result slot:} invalid-policy
KVM rows and failure evidence. \\
Lower-filesystem preservation & Lower filesystem keeps permissions, data path,
writes, persistence, and page-cache behavior for selected objects. & Broader
filesystems may own or interpose on data and write paths when their object
model requires it. & \textbf{Result slot:} permission/write/data preservation
rows tied to each admitted workload. \\
Source responsibility map & Policy responsibility is limited to lookup/readdir
selection and visibility for the admitted oracle. & The boundary account maps
source-system responsibilities such as create, unlink, rename, daemon/runtime
state, COW/checkpoint state, data path, or metadata persistence to the
filesystem methods or services a broader design would need to own. &
\textbf{Result slot:} per-source rows naming the owned methods, state
responsibility, privileged code surface, and policy code size with source
evidence. \\
\hline
\end{tabular}
\caption{RQ3 boundary result slots. The comparison is responsibility and
containment for the same oracle-relevant behavior, not a claim that custom
filesystems are less expressive.}
\label{tab:rq3-evidence}
\end{table}

For the Agent workspace matrix, the executed rows compare the complete source
lifecycle's name-resolution decisions with \namei and feature-equivalent FUSE.
AgentFS, BranchFS, YoloFS, and custom or stackable filesystems serve as boundary
exemplars unless a source oracle admits them as same-oracle executed rows. For
the environment/cache matrix, the boundary account compares the cache-state
name-resolution policy with the original environment mechanism,
feature-equivalent FUSE ownership, and any custom or stackable filesystem
responsibility needed to own cache metadata, stale/corrupt state, or data-path
behavior.

RQ3 is answered positively when the same-oracle rows pass, the
real attach path rejects malformed decisions visibly, and the boundary account
shows that broader filesystem mechanisms would own methods, daemon state,
storage semantics, or data-path behavior that \namei leaves to VFS and the
lower filesystem. RQ3 fails if the admitted source oracle requires that broader
ownership in the common case. The corresponding result slot reports ownership,
invalid-decision, and lower-filesystem preservation cells in
Table~\ref{tab:rq3-evidence}.

\subsection{Evaluation Scope And Limitations}

Rows that require final-file target selection, including environment/cache
file-object evidence, are outside the evidence reported here until that support
exists. Service/config is included only when a concrete lookup-time source
oracle exists and remains motivating and conditional scope otherwise. RQ
answers come from reviewed KVM source-oracle runs,
feature-equivalent FUSE comparisons, and ownership-boundary evidence.
Materialized namespace mechanisms, eBPF LSM hooks, and native production
mechanisms remain related work and background context unless a selected source
workload makes one of them part of the source behavior being evaluated.
